% =====================================================================
\section{Where fault-tolerant estimates stand}
\label{sec:ft}

\subsection{How to read a resource estimate}
A fault-tolerant quantum computer combines many \emph{physical} qubits into fewer \emph{logical} qubits that are protected against errors. A resource estimate for an algorithm reports some of the following, and the differences between them are large.
\begin{itemize}
\item \textbf{Logical width}: the logical qubits the algorithm occupies. Sources count it differently. Some give the minimum workspace; others include routing space or memory, or exclude the factories described next.
\item \textbf{Non-Clifford gate count}: the number of T gates, or of Toffoli (CCZ) gates. These are the expensive operations. Each consumes a specially prepared ``magic'' state produced by a \emph{factory}, an error-corrected subroutine that can occupy most of the machine. T and Toffoli counts are different units and are not interchangeable here.
\item \textbf{Physical qubits and runtime}: derived from a code, a physical error rate, a code-cycle time and a classical reaction time. Runtime may be quoted for one shot or as an expected value that includes repeats after detected failure.
\item \textbf{What is left out}: commonly the preparation of the initial state, the number of repetitions needed for a reliable answer, and every calculation after the first.
\end{itemize}
Fig.~\ref{fig:dumbbell} shows the five cases in which a modelled quantum time and an estimated classical time exist for one model and tolerance. None of the plotted values is a measured runtime, and the cases below explain what each one covers. The full list of resource records is Table~\ref{tab:worked} in Appendix~\ref{app:models}.

\begin{figure}[htbp]
\centering
\begin{tikzpicture}
\begin{axis}[
  xmode=log, width=0.56\textwidth, height=8.0cm,
  xmin=0.2, xmax=3e7, ymin=0.4, ymax=5.6, y dir=reverse,
  ytick={1,...,5},
  yticklabels={{Fermi--Hubbard 6$\times$6, half filling\\{\scriptsize\color{black!60} model schedule vs 4-core DMRG estimate}}, {Fermi--Hubbard 10$\times$10, half filling\\{\scriptsize\color{black!60} model schedule vs 4-core DMRG estimate}}, {$J_1$--$J_2$ Heisenberg 10$\times$10\\{\scriptsize\color{black!60} model schedule vs 4-core DMRG estimate}}, {FeMoco, 76 orbitals, 1 kcal/mol\\{\scriptsize\color{black!60} model expectation vs idealized estimate, not run}}, {P450 compound I, 58 orbitals\\{\scriptsize\color{black!60} model expectation vs 2022 source estimate}}},
  yticklabel style={font=\footnotesize, align=right},
  xlabel={\footnotesize Published time estimate (hours, log scale)},
  xticklabel style={font=\footnotesize},
  xmajorgrids, grid style={black!10}, axis y line*=left, axis x line*=bottom, clip=false,
  legend style={font=\scriptsize, at={(0.5,-0.15)}, anchor=north, draw=black!25, fill=white, legend columns=2}, legend cell align=left]
\draw[black!20, line width=1.4pt] (axis cs:3.4167,1) -- (axis cs:61.111,1);
\draw[black!20, line width=1.4pt] (axis cs:29.167,2) -- (axis cs:1.9722e6,2);
\draw[black!20, line width=1.4pt] (axis cs:15.667,3) -- (axis cs:21.944,3);
\draw[black!20, line width=1.4pt] (axis cs:1.11,5) -- (axis cs:720,5);
\draw[qcU!60, line width=1.4pt] (axis cs:5.26,4) -- (axis cs:42.08,4);

\addplot[only marks, mark=o, mark size=3pt, color=qcF, line width=1pt] table[x=x,y=y] {figdata/dumbbell_q_est.dat};
\addlegendentry{quantum, modelled schedule}
\addplot[only marks, mark=triangle, mark size=3.2pt, color=qcF, line width=0.9pt] table[x=x,y=y] {figdata/dumbbell_q_par.dat};
\addlegendentry{quantum, alternative schedule}
\addplot[only marks, mark=diamond, mark size=3.4pt, color=qcU, line width=1pt] table[x=x,y=y] {figdata/dumbbell_q_scen.dat};
\addlegendentry{quantum, repetition scenario (ours)}
\addplot[only marks, mark=square, mark size=2.8pt, color=black!65, line width=1pt] table[x=x,y=y] {figdata/dumbbell_c_est.dat};
\addlegendentry{classical, estimate}
% generated by ledger/make_figdata.py
\node[font=\scriptsize, inner sep=5pt, anchor=south] at (axis cs:61.1111,1) {61\,h};
\node[font=\scriptsize, inner sep=5pt, anchor=south] at (axis cs:1.97222e6,2) {$\sim$225\,y};
\node[font=\scriptsize, inner sep=5pt, anchor=south] at (axis cs:21.9444,3) {22\,h};
\node[font=\scriptsize, inner sep=5pt, anchor=north] at (axis cs:5.26,4) {5.3\,h};
\node[font=\scriptsize, inner sep=5pt, anchor=south] at (axis cs:22.9,4) {22.9\,h};
\node[font=\scriptsize, inner sep=5pt, anchor=north] at (axis cs:73.0556,5) {2022: 25--73\,h};
\node[font=\scriptsize, inner sep=5pt, anchor=south] at (axis cs:1.88,5) {2026: 1.1--1.9\,h};
\node[font=\scriptsize, inner sep=5pt, anchor=south] at (axis cs:720,5) {$\sim$1 month};

\node[font=\scriptsize, anchor=west, qcUd] at (axis cs:60,4.02) {\repAttempts{} or \repAttemptsSplit{} accepted runs: \repHours, \repHoursSplit\,h};
\end{axis}
\end{tikzpicture}
\caption[Modelled quantum and estimated classical times]{Modelled quantum times and estimated classical times for five cases in which both sides publish a value for one model and tolerance. No point is a measured runtime, and the type of each value is printed under its row. \emph{Lattice models:} circles are the single-factory phase-estimation schedules of the source's Table~S12 and triangles its 16-factory schedules, which use \parQubitLo--\parQubitHi{} times the physical qubits; squares are the source's DMRG estimates on four cores of one CPU for a variational state within 0.01 of the extrapolated energy~\cite{yoshioka2024}. \emph{FeMoco:} circle and triangle are the minimum-volume and fastest tabulated rows of the 2026 dense-code compilation~\cite{low2026dense}, each an expected duration that already includes repeats after detected logical failure; diamonds are our scenario for \repAttempts{} or \repAttemptsSplit{} accepted executions at a hypothetical hit probability of 0.5 (Section~\ref{sec:qpe}); the square is an idealized conversion for a DMRG calculation that was not performed~\cite{zhai2026}. \emph{P450:} circles are the 2022 estimates at physical error rates $10^{-3}$ and $10^{-5}$~\cite{goings2022}, triangles the 2026 rows~\cite{low2026dense}, and the square the 2022 source's own classical estimate. A connector shows that two values refer to one model. It implies no ranking where the two sides differ in confidence semantics or machine assumptions. Generated from the ledger.}
\label{fig:dumbbell}
\end{figure}


\subsection{Lattice models of correlated electrons}
\label{sec:lattice}
\textbf{Question.} The Fermi--Hubbard model, electrons hopping on a lattice with on-site repulsion $U$, and the frustrated $J_1$--$J_2$ Heisenberg spin model are the standard minimal models of correlated materials. The output considered here is a ground-state energy to a total-energy accuracy of 0.01 in the model's units.

\textbf{What was done.} Yoshioka et al.\ estimate the cost of quantum phase estimation in the surface code at a physical error rate of $10^{-3}$ and compare it with DMRG~\cite{yoshioka2024}. For the half-filled Hubbard model at $U/t=4$ on a $10\times10$ cylinder, the modelled schedule with one magic-state factory takes $1.05\times10^5$\,s (\hubTenQh{} hours) on $8.9\times10^5$ physical qubits; with sixteen factories it takes $1.46\times10^4$\,s on $5.4\times10^6$. For the Heisenberg model at $J_2/J_1=0.5$ on the same lattice the two schedules take $5.64\times10^4$\,s and $4.32\times10^3$\,s. The classical figures are the source's estimates of the DMRG time needed to produce a variational state within 0.01 of the extrapolated energy: $7.1\times10^9$\,s for the Hubbard case and $7.9\times10^4$\,s for the Heisenberg case (Fig.~\ref{fig:hubbard}).

\begin{figure}[htbp]
\centering
\begin{tikzpicture}
\begin{axis}[
  xmode=log, ymode=log, width=0.68\textwidth, height=6.6cm,
  xmin=12, xmax=130, ymax=3e10,
  xtick={16,32,36,64,100}, xticklabels={16,32,36,64,100},
  xlabel={\footnotesize Lattice sites (4$\times$4, 4$\times$8, 6$\times$6, 4$\times$16 and 8$\times$8, 10$\times$10)},
  ylabel={\footnotesize Estimated time to $\Delta E=0.01$ (s)}, ymin=1, clip=false,
  ticklabel style={font=\footnotesize}, grid=major, grid style={black!10},
  legend style={font=\scriptsize, at={(0.02,0.98)}, anchor=north west, draw=black!25, fill=white}, legend cell align=left]
\addplot[black!55, mark=square*, line width=0.9pt, mark size=2.2pt, dashed, mark options={solid}] table[x=x,y=y] {figdata/hubbard_U1.dat};
\addlegendentry{DMRG, square, $U/t=1$}
\addplot[qcF, mark=*, line width=1.2pt, mark size=2.4pt] table[x=x,y=y] {figdata/hubbard_U4.dat};
\addlegendentry{DMRG, square, $U/t=4$}
\addplot[qcU, mark=triangle*, line width=1.0pt, mark size=2.8pt, dashdotted, mark options={solid}] table[x=x,y=y] {figdata/hubbard_U8.dat};
\addlegendentry{DMRG, square, $U/t=8$}
\addplot[only marks, mark=o, mark size=2.6pt, qcF, line width=0.8pt] table[x=x,y=y] {figdata/hubbard_strip_U4.dat};
\addlegendentry{strips $4\times8$, $4\times16$ (hollow, same shapes)}
\addplot[only marks, mark=square, mark size=2.2pt, black!55, line width=0.8pt, forget plot] table[x=x,y=y] {figdata/hubbard_strip_U1.dat};
\addplot[only marks, mark=triangle, mark size=2.8pt, qcU, line width=0.8pt, forget plot] table[x=x,y=y] {figdata/hubbard_strip_U8.dat};
\addplot[only marks, mark=diamond*, mark size=3.6pt, qcQ] table[x=x,y=y] {figdata/hubbard_qpe_s.dat};
\addlegendentry{QPE, 1 factory (modelled)}
\addplot[only marks, mark=diamond, mark size=3.6pt, qcQ, line width=0.9pt] table[x=x,y=y] {figdata/hubbard_qpe_p.dat};
\addlegendentry{QPE, 16 factories (modelled)}
\end{axis}
\end{tikzpicture}
\caption[Half-filled Hubbard model: classical estimate and quantum model]{Half-filled two-dimensional Fermi--Hubbard model on cylinders. Filled symbols joined by lines (squares and a dashed line for $U/t=1$, circles and a solid line for $U/t=4$, triangles and a dash-dotted line for $U/t=8$): the source's estimated DMRG time, on four cores of one CPU, to a variational state within 0.01 of the extrapolated total energy, for square lattices (its Table~S2). Hollow symbols: narrow strips, which are far cheaper at equal site count and are not joined to the squares. Diamonds, filled and hollow: the source's modelled phase-estimation schedules at $U/t=4$ with one and with sixteen magic-state factories (its Table~S12)~\cite{yoshioka2024}. This instance has no sign problem for determinant and auxiliary-field quantum Monte Carlo~\cite{hirsch1985}; no Monte Carlo cost is shown because none has been compiled at this tolerance. Generated from the ledger.}
\label{fig:hubbard}
\end{figure}


\textbf{Evidence.} Both sides are models. Against this baseline the single-factory quantum schedule is faster by a factor of \ratioHubTen{} for the Hubbard lattice and \ratioHeis{} for the Heisenberg lattice. Three properties of the baseline limit what those ratios mean.
\begin{itemize}
\item \emph{Machine.} The DMRG times are for four cores of one CPU. The Hubbard figure is \hubTenCoreH{} core-hours, about three times the classical FeMoco campaign of Section~\ref{sec:chem} and the allocation equivalent of roughly \hubTenFugakuH{} hour of all Fugaku cores. That equivalence is a statement of budget and not a runtime, since DMRG does not parallelize ideally; the source itself expects about a tenfold gain from GPUs or several nodes, and GPU implementations now reach bond dimensions above $10^4$ on one node~\cite{legeza2026}.
\item \emph{Criterion.} The classical time is for a variational \emph{state} of the target accuracy. Phase estimation returns an energy, and the source notes that estimating the energy alone by extrapolation is much faster.
\item \emph{Method.} At half filling with nearest-neighbour hopping on a bipartite lattice, the Hubbard model has no sign problem for determinant and auxiliary-field quantum Monte Carlo~\cite{hirsch1985,leblanc2015}. The sign problem is the statistical cancellation that makes Monte Carlo expensive for most fermion problems; here it is absent. No Monte Carlo time-to-solution has been compiled at this tolerance, so the best classical baseline for this instance is unknown. For the Heisenberg benchmark on a periodic cluster, a transformer neural quantum state reaches a total energy 0.21 below a DMRG value at bond dimension 8{,}192~\cite{rende2024}. That is another geometry, but it shows that the cited DMRG calculation is not the best variational energy in that tabulated benchmark.
\end{itemize}

\textbf{Unresolved.} A measured classical speedup above \thrHeis{} over the four-core estimate would reverse the single-factory Heisenberg comparison, and one above \thrHeisPar{} would reverse the sixteen-factory comparison; neither has been measured. The Heisenberg case is therefore a small modelled margin whose sign depends on implementation and preparation assumptions. The Hubbard margin is large against DMRG and untested against Monte Carlo. On the quantum side, adiabatic state preparation is assumed to be cheap on the basis of fits up to 48 qubits. Doping, frustration or longer-range hopping restore the sign problem and define separate candidates, for which neural, tensor-network and embedding methods are the classical side (class 36). A recent experiment measured pairing correlations in a doped Hubbard model on trapped-ion hardware~\cite{granet2025}; pairing is a different observable from the ground-state energy.

\textbf{Next test.} A Monte Carlo, neural-state or parallel DMRG energy on the same cylinder at the same tolerance with a measured time; and a quantum estimate that includes preparation.

\subsection{Transition-metal chemistry: FeMoco and P450}
\label{sec:chem}

\textbf{Question.} FeMoco is the iron--molybdenum cofactor at which the enzyme nitrogenase reduces nitrogen, and cytochrome P450 is a family of enzymes that metabolize drugs. Their reactivity depends on the energies of closely spaced electronic states. Calculations treat an \emph{active space}: a chosen set of electrons and orbitals handled exactly, here 76 orbitals for FeMoco and 58 for the P450 intermediate compound~I. An active-space energy is not the chemistry of the enzyme. The usual target is \emph{chemical accuracy}, about 1\,kcal/mol or 1.6\,mHa.

\textbf{Quantum estimates.} Quantum phase estimation extracts an energy from a prepared state. Its cost for these models has fallen steadily. Tensor hypercontraction (2021) needs 2{,}196 logical qubits and $3.2\times10^{10}$ Toffoli gates for the 76-orbital FeMoco model; the frequently quoted four million physical qubits and four days belong to the smaller 54-orbital model, at 2{,}142 logical qubits and $5.3\times10^9$ Toffoli gates~\cite{lee2021thc}. A 2026 compilation into a dense surface-code layout tabulates a frontier of physical qubits against runtime for both molecules (Fig.~\ref{fig:dc76})~\cite{low2026dense}. For FeMoco it gives \femoT{} expected hours at \femoSevMinvolQ{} million physical qubits with \femoTofMinvol{} Toffoli gates (the row of minimum space--time volume), and \femoFastT{} hours at \femoSevFastQ{} million with \femoTofFast{} (the fastest row tabulated). For P450 it gives \pfMinvolH{} hours at 1.03 million and \pfFastH{} hours at 2.18 million, against 73 hours on 4.6 million in the 2022 estimate for the same model~\cite{goings2022}: about \pfSpeed{} times less time on \pfQubits{} times fewer qubits at the same error rate.

\begin{figure}[htbp]
\centering
\begin{tikzpicture}
\begin{groupplot}[group style={group size=2 by 1, horizontal sep=1.9cm}, width=0.47\textwidth, height=6.0cm, ticklabel style={font=\footnotesize}, label style={font=\footnotesize}, title style={font=\footnotesize}, xmode=log, ymode=log, xmin=5e4, xmax=8e6, grid=major, grid style={black!10}, xlabel={Physical qubits}, legend style={font=\scriptsize, draw=black!25, fill=white}, legend cell align=left]
\nextgroupplot[ymin=1, ymax=1e5, ylabel={Expected accepted runtime (hours)}, title={FeMoco, 76 orbitals}, legend pos=north east]
\addplot[qcF, mark=*, mark size=2.6pt, line width=0.6pt, densely dotted, mark options={solid}] table[x=x,y=y] {figdata/femoco76_dense.dat};
\addlegendentry{dense code, 2026}
\nextgroupplot[ymin=0.5, ymax=1e5, title={P450 Cpd I, model X (58 orbitals)}, legend pos=north east]
\addplot[qcF, mark=*, mark size=2.6pt, line width=0.6pt, densely dotted, mark options={solid}] table[x=x,y=y] {figdata/p450dc_dense.dat};
\addlegendentry{dense code, 2026}
\addplot[only marks, mark=square*, mark size=2.6pt, black!55] table[x=x,y=y] {figdata/p450dc_prior.dat};
\addlegendentry{THC, 2022}
\end{groupplot}
\end{tikzpicture}
\caption[FeMoco and P450: selected rows of the 2026 quantum cost frontier]{Each curve trades physical qubits against modelled time for one energy calculation; state preparation and the reliability of a useful answer are separate. Selected rows of the space--time frontier for phase estimation on the 76-orbital FeMoco model and the 58-orbital P450 model~\cite{low2026dense}. The vertical axis is the source's expected runtime. Dotted lines are guides to the eye between selected rows. Circles: rows of the 2026 dense-code compilation. Squares: the 2022 estimates for the same P450 model~\cite{goings2022}. Row selection and assumptions are in Appendix~\ref{app:models}. Generated from \texttt{resources.csv}.}
\label{fig:dc76}
\end{figure}


Two features of these numbers matter for any comparison. First, each runtime is an \emph{expected} value: the source divides the one-shot runtime by one minus the probability of a detected logical failure. For the FeMoco minimum-volume row that is \femoShot{} hours per shot at a failure probability of \femoPfail. Toward the low-qubit end failure dominates: at \femoSevLowQ{} physical qubits a 529-hour shot fails with probability \femoLowPfail{} and the expected time is 716 hours. Second, a successful shot is not yet a useful answer. Phase estimation returns the ground-state energy only with a probability set by how much the prepared state overlaps the ground state, so several accepted executions are needed, and no overlap has been measured for these models. Section~\ref{sec:qpe} works through the consequence: at a hypothetical hit probability of 0.5, \repAttempts{} or \repAttemptsSplit{} accepted executions, depending on how the failure budget is divided, or \repHours{} to \repHoursSplit{} expected hours.

\textbf{Classical results.} Zhai et al.\ computed the ground-state energy of the 76-orbital FeMoco model with unrestricted DMRG at bond dimensions up to 18{,}000 and high-order coupled cluster, starting from 78{,}750 broken-symmetry guesses~\cite{zhai2026}. The result agrees between methods to $\pm0.31$\,kcal/mol and has a conservative range of $\pm1.25$\,kcal/mol. Both rest on extrapolation, so this is approximately chemical accuracy with an empirical uncertainty, not a certified interval. The insight behind it is structural: the hard part is ranking many candidate states that are each close to a single configuration. The source's cost statements are three separate records.
\begin{itemize}
\item \emph{Measured:} the DMRG campaign took 2{,}768{,}994 core-hours, of which one sweep at bond dimension 18{,}000 took 317{,}440.
\item \emph{Modelled, not performed:} two variational sweeps at a bond dimension near 393{,}000 would take \femoClassCoreH{} core-hours by the fitted cost law, which the source converts to \femoClassIdeal{} hours on an exascale machine at ideal efficiency. The same conversion gives \femoUcc{} seconds for a coupled-cluster estimate. The source describes these as upper bounds that omit practical implementation factors.
\item \emph{Not costed:} constructing and ranking the candidate states.
\end{itemize}
For the P450 model the 2022 source measured DMRG on a smaller model and \emph{estimated} 36{,}564 CPU-hours, about a month, for model X at bond dimension 3{,}000, with no guarantee that this bond dimension suffices~\cite{goings2022}.

\textbf{What the comparison shows.} Three questions should be kept apart. \emph{Can classical methods reach the FeMoco model energy?} Yes, to an extrapolation-dependent uncertainty; this is well supported, and it removes this active-space model, at the source's extrapolation-dependent energy target, as an example of classical intractability; it certifies no error interval and says nothing about other observables or larger models. \emph{Which workflow reaches a stated target fastest?} Not settled. The \femoT-hour quantum figure is a model expectation for one accepted execution from a supplied state. The \femoClassIdeal-hour classical figure is an idealized conversion for a calculation that was not run. Neither includes the search for candidate states, which both approaches need. \emph{Which gives the required reliability?} Not settled: one side offers a per-execution probability that depends on an unmeasured overlap, the other an empirical extrapolation. For P450 the scientific target is a gap between spin states, which needs at least two calculations on either side with errors that must be shown to cancel; a one-hour quantum runtime for a single energy is not yet a comparison on that task.

\textbf{Related records.} A model of P450 with about 500 orbitals would need 9{,}000 logical qubits and $1.5\times10^{12}$ Toffoli gates~\cite{goings2022}. A ruthenium catalyst for carbon fixation was estimated in 2022 at 2{,}740 logical qubits and $5.4\times10^{11}$ T gates, a month to 260 years depending on the hardware scenario~\cite{beverland2022}; the 2026 compilation gives 0.46 hours at 1.57 million physical qubits for a different, 56-orbital model of the same catalyst family~\cite{low2026dense}, so the older figure is a historical example and not the current frontier. Across chemistry more generally, Lee et al.\ find no evidence of an exponential advantage in the classes they studied, which leaves polynomial advantages and specific strongly correlated cases open~\cite{lee2023}.

\textbf{Next test.} A measured or bounded overlap for the prepared state, with the full protocol costed, including candidate search and, for P450, both spin states; and a classical result at matched confidence for a model larger than the one already solved.

\subsection{Core-level spectroscopy}
\label{sec:xas}
\textbf{Question.} X-ray absorption spectra reveal oxidation states in battery cathodes, and interpreting them needs calculated spectra of many excited states. The output is a spectrum at a stated broadening.

\textbf{What was done.} Fomichev et al.\ give a time-domain quantum algorithm for a Li$_4$Mn$_2$O cluster with 22 electrons in 18 orbitals~\cite{fomichev2025}. In the current version of their resource table, the longest circuit uses \xasCirc{} Toffoli gates and the whole sampled spectrum \xasAlg{}, about \xasCircuits{} times as many, for a budget of 2{,}500 shots from an initial state built from $10^4$ Slater determinants. The quoted 100 logical qubits are the minimum workspace, excluding factories; the timed layout uses 350. On that layout the source counts $1.58\times10^{11}$ logical cycles for the whole algorithm and describes this as less than two days at megahertz logical clock rates; at exactly 1\,MHz it is \xasHours{} hours. The previous version of the paper gave \xasAlgOld{} gates and \xasHoursOld{} hours: the revision raised the gate count by a factor of \xasGateRatio{} and the cycle count by \xasCycleRatio.

\textbf{Evidence.} Among the chemistry and materials estimates of this section, this is the only one that costs its protocol end to end at the logical level: preparation, sampling and the full spectrum. Its assumptions are not validated. No machine has demonstrated a megahertz logical clock. The classical comparison is the source's estimate of about 32 days for 1{,}000 excited states by a restricted-active-space method, 0.76 hours per state, which the source itself calls likely optimistic on the classical side. No classical spectrum for this cluster is given; the algorithm is checked against an exact spectrum only for a system small enough to simulate.

\textbf{Unresolved and next test.} How well the $10^4$-determinant state represents the true initial state; how secure the perturbative estimate of spectral error is at this size; and whether the restricted-active-space estimate represents the best classical spectrum method, including neural excited-state solvers~\cite{pfau2024}. A classical spectrum for the same cluster at decision accuracy would answer the last.

\subsection{Linear systems and fluid dynamics: a negative case}
\label{sec:linear}
Sparse linear systems are the workhorse of engineering simulation. Finite-element structural analysis, implicit computational fluid dynamics (CFD), electromagnetics and weather prediction all spend most of their time solving $Ax=b$ for matrices with millions to trillions of unknowns. Quantum linear solvers are regularly proposed as accelerators for this work. Theory and the published estimates do not establish that they are. The evidence is of four different kinds, and we keep them apart.

\textbf{What a quantum linear solver does.} The algorithm of Harrow, Hassidim and Lloyd (HHL) prepares a quantum state whose amplitudes are proportional to the solution $x=A^{-1}b$~\cite{hhl2009}. For a sparse $N\times N$ matrix its cost grows as a polynomial in $\log N$, in the condition number $\kappa$ (the ratio of the largest to the smallest singular value, a measure of how strongly an input error is amplified in the solution; the value that matters is that of the matrix as encoded, after any preconditioning) and in the inverse accuracy; later solvers reach a cost linear in $\kappa$ and logarithmic in the inverse accuracy of the state, which is optimal~\cite{costa2022}. The general problem is BQP-complete under its stated access and output model, so a general efficient classical algorithm would imply efficient classical simulation of quantum computation~\cite{hhl2009}. This supports a conditional complexity advantage. It is not a proved exponential separation, and it says nothing by itself about the engineering instances considered here. The attractive regime is defined by four conditions~\cite{aaronson2015}: the right-hand side can be loaded as a quantum state quickly; the matrix entries can be computed on demand or the matrix otherwise encoded in a circuit (a block encoding, which embeds a scaled copy of the matrix in a larger unitary and whose scale factor and cost can cancel a nominal gain); $\kappa$ grows at most polylogarithmically in $N$; and the user needs only a few numbers derived from $x$, not $x$ itself. Table~\ref{tab:linreq} compares each condition with what engineering problems supply. The conditions mark where an advantage is possible; they do not guarantee one, since structured instances can be easy for both kinds of machine.

\begin{table}[htbp]
\centering\small
\caption{Conditions for an exponential speedup from a quantum linear solver, and what discretized engineering problems supply.}\label{tab:linreq}
\begin{tabular}{@{}L{3.0cm} L{6.0cm} L{6.8cm}@{}}
\toprule
Condition & In finite-element and CFD problems & Consequence\\\midrule
Small condition number & For a second-order elliptic operator on a uniform mesh in $d$ dimensions, $\kappa$ grows as $N^{2/d}$; practical CFD matrices reach $10^6$ and more~\cite{goldfriend2026} & Quantum cost cannot be made sublinear in $\kappa$ under standard complexity assumptions~\cite{hhl2009}. Classical multigrid solves the same systems in time proportional to $N$, independent of $\kappa$~\cite{brandt1977}\\
Size independent of accuracy & $N$ is set by the accuracy wanted: halving the discretization error multiplies $N$ by a fixed power & A speedup in $N$ becomes a speedup in the inverse accuracy, which is at most polynomial in fixed dimension~\cite{montanaro2016}\\
Efficient input & Meshes, geometries and boundary data are classical; the matrix must be block-encoded, which inflates the effective condition number~\cite{lapworth2025} & Encoding cost and classical preprocessing enter the total~\cite{lapworth2022}\\
Few numbers out & A drag coefficient is one number, estimated to accuracy $\epsilon$ at a cost growing as $1/\epsilon$; a pressure or velocity field is $N$ numbers & Scalar outputs lose the logarithmic accuracy dependence of classical iteration; full fields need at least $N$ state preparations\\
Unitary-like dynamics & Advection--diffusion operators are dissipative and non-normal & Proven overheads for quantum differential-equation solvers from both properties~\cite{an2025}\\
Weak nonlinearity & Turbulent flow is strongly nonlinear and chaotic & Linearization is efficient only for weak nonlinearity~\cite{liu2021c}; no efficient quantum algorithm for chaotic regimes under stated assumptions~\cite{lewis2024}\\
\bottomrule
\end{tabular}
\end{table}

\textbf{First kind of evidence: theorems for a stated problem and access model.} Montanaro and Pallister analyzed the finite-element method with the quantum solver and a classical conjugate-gradient solver held to the same final accuracy, for a linear functional of the solution of an elliptic boundary-value problem. Earlier work had claimed an exponential speedup. With accuracy accounted for, the speedup is at most polynomial when the spatial dimension is fixed and the solution is smooth; for a problem in three space dimensions and time it is roughly quadratic at best, and in low dimension with a smooth solution the quantum algorithm can be slower~\cite{montanaro2016}. With optimal preconditioning on both sides, their complexity table has the classical cost growing as $\epsilon^{-d/2}$ in $d$ dimensions and the quantum cost as $\epsilon^{-1}$, up to logarithms and norms of the solution: the two are equal at $d=2$ and the gap is quadratic only at $d=4$. Linden, Montanaro and Shao compared five classical and five quantum algorithms for the heat equation. The approach based on a quantum linear solver is never faster than the best classical method, and the best quantum method, amplitude estimation applied to a classical random walk, gains at most a quadratic factor in the accuracy~\cite{linden2022}. For time-dependent problems, An et al.\ prove worst-case lower bounds showing that quantum solvers pay for two kinds of departure from quantum dynamics, dissipation and non-normality, and conclude that quantum algorithms work best for quantum dynamics~\cite{an2025}. For nonlinear problems, the Carleman linearization that underlies most proposals is efficient only when nonlinearity is weak relative to dissipation, and the problem is intractable in the worst case beyond a threshold~\cite{liu2021c}; Lewis et al.\ show that, for the evolution and output encoding of their theorem, no efficient quantum algorithm recovers the state of a chaotic system over long times~\cite{lewis2024}. A long-horizon trajectory, a short-time prediction, a time-averaged drag and a stationary distribution are different computational products, and these results bear on them differently. What these results leave, for the problems they cover, is a polynomial speedup; Section~\ref{sec:breakeven} illustrates how sensitive a low-degree polynomial speedup is to the speed of error-corrected hardware.

\textbf{Second and third kinds: estimates for stated algorithms, and conditional proposals.} Table~\ref{tab:linstudies} lists seven studies. Scherer et al.\ and Penuel et al.\ each cost one algorithm for an application that had been put forward as promising: the electromagnetic-scattering problem for which an exponential speedup had been claimed~\cite{clader2013}, and drag calculation for ship-hull design. Their size is evidence about those algorithms and not a lower bound on every method. Jennings, Zhuang and Tu report conditional gains, each under an assumption named in the table, and the two preconditioning studies show what a full account must include. We read each source in the full text for the points quoted. None of the numbers enters a figure: they are not matched pairs on one accuracy.

\begin{table}[htbp]
\centering\small
\caption{Selected assessments of engineering quantum solvers: scope and unresolved costs. Each row gives one reported result and the main element missing from a comparison. Counts, assumptions and locators are in Table~\ref{tab:worked} and the ledger.}\label{tab:linstudies}
\begin{tabular}{@{}L{3.4cm} L{6.6cm} L{5.8cm}@{}}
\toprule
Study and problem & What it reports & What is missing\\\midrule
Scherer et al.~\cite{scherer2017}: electromagnetic scattering, $3.3\times10^8$ unknowns & An estimate for one algorithm: circuit depth $1.8\times10^{29}$ with oracles & Nothing that would rescue this compilation; it is not a bound on other methods\\
Penuel et al.~\cite{penuel2025}: drag on a sphere and ship hulls & A partial, optimistic estimate for one algorithm: qubits times T gates from $10^{21}$ to $10^{39}$, stated as lower bounds & Initial-state and boundary-condition costs; a time comparison (the product is a resource proxy); reconciliation with its public records\\
Jennings et al.~\cite{jennings2025}: incompressible lattice-Boltzmann flow & A repaired algorithm with bounded gain: at best $\mathrm{Re}^{3/4}$ for drag in three dimensions, and $\mathrm{Re}^{0.287}$ in two dimensions once its numerical conditioning is used & Evidence that any gain survives at the accuracy and Reynolds number of an engineering output\\
Zhuang et al.~\cite{zhuang2025}: Navier--Stokes on a $2^{80}$-point mesh & The only application-specific crossover: 42.6 days on 8.71 million physical qubits against a modelled 130 years & A classical solver given the same spectrally compact input and output\\
Lapworth and S\"underhauf~\cite{lapworth2025}: industrial CFD matrices & Preconditioning trades conditioning against encoding cost; in a worked example the total rises & An instance where the total falls, with the preprocessing credited to both sides\\
Goldfriend et al.~\cite{goldfriend2026}: hybrid nozzle-flow solver, $\kappa\approx80$ & Approximate coefficient preparation cuts one resource from 254 angles to 12, at twice the outer iterations & Total cost, success overhead and industrial conditioning\\
Tu et al.~\cite{tu2025}: generic sparse systems & A crossing in a kernel and hardware model from $N\approx2^{33}$ & Input and output costs; hardware with a 10\,ns logical gate\\
\bottomrule
\end{tabular}
\end{table}

\textbf{Fourth kind: what has not been shown.} None of these studies establishes a practical advantage over a matched, strong classical workflow at complete cost. The two compilations are tens of orders of magnitude from feasibility. The drag study makes the distance concrete: the classical runs it takes as its threshold of utility finished in hours on one workstation, with a turbulence model in place of resolved turbulence, while the resource products it derives as lower bounds start at $10^{21}$. The source sets a fitted exponent of 2.68 for that product beside a classical cost growing as $\mathrm{Re}^3$; the product is a resource proxy, and the two exponents are not a ratio of wall times. Its public per-instance records reach a product of $7.9\times10^{40}$ at $\mathrm{Re}=10^8$, 79 times the upper end of the published range, and the selection behind that range is not documented~\cite{qsub2025}; we quote the published range and record the difference as unresolved. The one application-specific crossover is reported for a mesh of $2^{80}$ points. Its quantum side relies on a spectrally compact input and output. Its classical side is an operation count for banded Cholesky factorization of the full system of dimension $2^{82}$, divided by a benchmark rate, and uses none of that compactness~\cite{zhuang2025}. The informative comparison is with a classical spectral solver given the same compact representation, and it has not been made; Section~\ref{sec:opps} sets it out as immediate future work. The crossing in the kernel model assumes a logical gate every 10\,ns, a hundred times the megahertz clock that Section~\ref{sec:xas} already notes is undemonstrated, and excludes input and output~\cite{tu2025}. The proposal that suggested an exponential advantage for fluid simulation by the lattice-Boltzmann route~\cite{li2025fluid} belongs to the family whose bottlenecks the two later studies document. A 2026 review reaches the same qualitative conclusion: any advantage is problem-dependent and governed by condition number, state preparation and measurement~\cite{cecile2026}. Solver research continues, including algorithms whose cost depends on an effective condition number of the instance and not on the worst case~\cite{dalzell2026}, but these do not remove the input and output conditions.

\textbf{The classical side keeps moving.} The incumbent here is strong and improving. For operators with a working multigrid hierarchy, multigrid-preconditioned Krylov solvers scale to the largest machines with cost near linear in the number of unknowns; nonsymmetric and nonnormal systems need other preconditioners and are harder. The fair baseline is a solver tuned to the instance. Learned models now add to them: machine-learned closures allow coarser grids at fixed accuracy~\cite{kochkov2021}, learned emulators produce medium-range weather forecasts in minutes on one accelerator~\cite{lam2023}, and many flow fields have compact tensor-network descriptions~\cite{mello2025}. A quantum method must beat this moving baseline, not the classical algorithm of the original speedup proof.

\textbf{Where a quantum linear solver could still matter.} The four conditions are most naturally met when the matrix comes from quantum physics or has structure a circuit can generate. These are the natural candidates, not the only ones the evidence permits. Response functions of a many-electron system require the inverse of a shifted Hamiltonian applied to a physical state: the matrix is encoded natively, the state is prepared by the simulation itself, and the wanted output is a scalar. High-dimensional linear problems are a second case. Babbush et al.\ give an exponential speedup for simulating $2^n$ coupled classical oscillators whose masses and couplings can be computed on demand, with a kinetic energy as output, and show the problem is as hard as general quantum computation~\cite{babbush2023osc}. That instance is not a mesh loaded from a design file, and for diffusion-type equations classical Monte Carlo already avoids the growth of cost with dimension. A third case is a linear solve inside a quantum workflow that must run anyway.

\textbf{Assessment.} For the engineering linear-system and flow problems reviewed, no practical quantum advantage over a matched, strong classical workflow is established. On present evidence we judge quantum linear solvers unlikely to accelerate the sparse solves and time stepping that dominate engineering simulation as it is practised. That is our assessment of this literature and not a theorem: the limits above hold for stated problems, algorithm families and compilations, and structured inputs and selected observables leave conditional polynomial opportunities whose complete costs and strongest classical competitors are unresolved. The test that would change the assessment is an end-to-end estimate, with input loading, readout and constant factors, that beats a tuned multigrid or Monte Carlo solver at an industrial condition number, for an output an engineer would use.

\subsection{A contrast: real-time gauge theory}
\label{sec:gauge}
Lattice gauge theory shows the opposite situation. For real-time evolution and for matter at finite density, standard importance-sampling methods face severe sign or phase problems: cancellations make Monte Carlo estimates statistically expensive. No general classical method with established accuracy and cost is known for the three-dimensional target, although restricted regimes and approximate methods still supply necessary baselines. What is missing is a controlled quantum estimate. Published T-gate counts for one fiducial calculation, the shear viscosity of SU(3) gauge theory on a $10^3$ lattice, range from about $10^9$ to more than $10^{48}$ across formulations and truncations (Fig.~\ref{fig:lgt} in Appendix~\ref{app:models})~\cite{lamm2026,ciavarella2025,rhodes2024}. The low values come from aggressive truncations that have not been shown to support a continuum limit, and all exclude state preparation and measurement. For the $\Sigma(72\times3)$ subgroup the primary gives $3.5\times10^{12}$ T gates for the Kogut--Susskind Hamiltonian and $7.1\times10^{12}$ for an improved one, on 27{,}000 data qubits before ancillas~\cite{osorio2025}; a count of $10^5$ to $10^6$ T gates is for one plaquette oracle and not for a simulation~\cite{draper2026}. These numbers are a strong motivation for further work and not a progress series at a common accuracy. Euclidean lattice QCD at zero density, by contrast, is mature classical practice; a quantum sampling algorithm has been demonstrated on a small model~\cite{yamamoto2022}, with no evidence yet of a practical gain over modern hybrid Monte Carlo.
